\scanpage{35}
% letterhead
\begin{center}
INSTITUT\\
DES HAUTES ÉTUDES SCIENTIFIQUES\\
BURES-SUR-YVETTE (S.-\&-O.)\\
928-48-53
\end{center}

Les \unsure{éléments|élémt} de $DR(M)$ sont les éléments de
$DR(M)_K\simeq V\otimes_{\mathbb Q}K$ de la forme $z=x+iy$ $(x,y\in V_{\mathbb R})$
avec $x=fx$, $y=-fy$ (on exprime $\tau z=z$ i.e. $f\sigma z=z$ i.e. $\sigma z=fz$)
\deleted{\unsure{C'est d.|Cela}} Donc l'espace $DR(M)$ admet une \inserted{$\mathbb Q$-structure de} façon naturelle,
\deleted{\unsure{de la structure}} on prend le sous-espace des
$x+iy$, $x,y\in V$, $x=fx$, $y=-fy$. Soit $W$ cet espace, il
mérite le nom de \uline{sous-espace de Betti \inserted{tordu}} de $DR(M)$.

\diagram{\begin{tikzcd}[ampersand replacement=\&,column sep=small,row sep=small]
H^*(X(K),\mathbb Q)=V \arrow[r,phantom,"\subset"] \arrow[d,phantom,"\cup"] \& DR(M)_K\simeq DR(M_K) \arrow[d,phantom,"\cup"]\\
H^*(X(\mathbb R),\mathbb Q)=V\cap W=V^f=W^{F} \arrow[r,phantom,"\subset"] \& W
\end{tikzcd}}

Il est stable par $f$. Mais
la formation de ce sous-espace
\unsure{ne|se} se comporte-t-il comme
il faut par $\otimes$ ? \uline{Oui}.

On a ainsi un deuxième foncteur fibre sur $M$
à valeurs sur $\mathbb Q$, mais il est isom \inserted{comme foncteur $\otimes$ \unsure{cohérent}} au
premier, \deleted{\unsure{une fois,}} comme on voit en écrivant
$V\simeq V^++V^-$, $W=W^++W^-$, et utilisant les
isomorphismes $V^+\simeq W^+$, $V^-\simeq W^-$ canoniques.
Cet isom $V\unsure{-|\to}W$ est \deleted{\unsure{dé.|f.}} induit par l'autom
de $DR(M_K)\simeq V_K\simeq W_K$ qui est l'identité sur
$V_K^+$, et la multiplication par $i$ sur $V_K^-$ ; Cependant, cet
isom \uline{n'est} \uline{pas} compatible avec le produit \uline{tensoriel}

Il semble que $\mathrm{Betti}\mathrm T$ et Betti ne sont \deleted{\unsure{pas}} peut être pas isom,
l'un est déduit \inserted{Betti \unsure{tordu}} de l'autre en tordant par le cocycle
défini par $f_{\mathbb R}$ ----
% archive sheet number: 34
